\setlength{\absparsep}{18pt}
\begin{resumo}
O aprendizado prático de Sistemas Operacionais (SO), em particular no ambiente Linux, constitui um pilar fundamental na formação de profissionais em Tecnologia da Informação. Tradicionalmente, as aulas práticas dependem do provisionamento de Máquinas Virtuais (VMs) locais ou de configurações em \textit{dual boot}. Entretanto, esse modelo impõe barreiras severas: elevado consumo de memória RAM e CPU, frequentes incompatibilidades de virtualização por \textit{hardware} (como desativação de VT-x na BIOS dos computadores pessoais dos alunos), tempo letivo substancial consumido pelo docente com suporte técnico de infraestrutura e reduzido engajamento discente frente a roteiros estáticos. Além disso, plataformas de simulação puramente executadas no cliente (\textit{client-side}) revelam-se inadequadas para avaliações somativas formais, uma vez que critérios de correção podem ser facilmente inspecionados ou adulterados via ferramentas de desenvolvedor (\textit{DevTools}). Este Trabalho de Conclusão de Curso apresenta a concepção e o desenvolvimento de uma plataforma educacional integrada para o ensino e avaliação prática de Sistemas Operacionais. A solução combina um núcleo de simulação Linux/POSIX de alta fidelidade desenvolvido em TypeScript puro com Sistema de Arquivos Virtual (\textit{Virtual File System} - VFS) executado de forma interativa no navegador, aliado a uma arquitetura de \textit{back-end} segura e desacoplada. O servidor é responsável pela gestão de turmas, controle temporal de exames mediante \textit{tokens} dinâmicos de liberação para aplicação presencial, agendamento de repescagens e segundas chamadas com autorizações restritas, e um motor de correção remota baseado em \textit{snapshots} assinados do estado do VFS. A validação técnica inicial conta com uma suíte de mais de 750 testes automatizados que cobrem chamadas POSIX e cenários reais de certificação. Conclui-se que a abordagem elimina a fricção de configuração local, otimiza o tempo didático em sala de aula e assegura a integridade pedagógica e avaliativa na disciplina de Sistemas Operacionais do curso de Tecnologia em Sistemas para Internet da UTFPR Câmpus Guarapuava.

\textbf{Palavras-chave:} Sistemas Operacionais. Linux. Simulação Web. Virtual File System. Educação em Computação. Avaliação Prática. UTFPR.
\end{resumo}
